%% Extraído de Textual/Textuais/Metodologia/05_Placas.tex (figura: fila de trabalho)
\begin{tikzpicture}[
        node distance=0.9cm and 1cm,
        font=\small,
        box/.style={draw, rounded corners, align=center, minimum height=0.9cm, minimum width=2.2cm},
        arrow/.style={-{Latex}, thick}
    ]
        \node[box] (hw) {Periférico\\(ADC, HX711, \dots)};
        \node[box, right=of hw] (thr) {\textit{Thread}\\bloqueante};
        \node[box, right=of thr] (msgq) {Fila de dados\\(\texttt{k\_msgq})};
        \node[box, right=of msgq] (wq) {Fila de trabalho\\única (\texttt{k\_work\_q})};
        \node[box, above right=0.5cm and 1cm of wq] (mgrs) {Gerenciadores\\de domínio};
        \node[box, below right=0.5cm and 1cm of wq] (nci) {NCI\\(rede UDP)};
        \node[box, right=of nci] (pc) {Aplicação\\de PC};
        \node[box, below=of msgq] (cmdq) {Fila de\\comandos};

        \draw[arrow] (hw) -- (thr);
        \draw[arrow] (thr) -- (msgq);
        \draw[arrow] (msgq) -- (wq);
        \draw[arrow] (wq) -- (mgrs);
        \draw[arrow] (wq) -- (nci);
        \draw[arrow] (nci) -- (pc);
        \draw[arrow] (pc) -- (nci);
        \draw[arrow] (nci) |- (cmdq);
        \draw[arrow] (cmdq) -- (wq);
    \end{tikzpicture}
